

The distribution $ \P_{X} $ of the random variable $ X$ is defined on the probability distribution
$ \P$ defined in part 1 of this exercise.

We would like to show that $ \P_{X} $ is a discrete distribution and its probability mass function is 
$ \sum_{n \in \Z} p_{n} p_{x-n}  $.

Starting with the definition of a distribution of a random variable, take $ A \subset E = \Z$,

\begin{align*}
    P_{X} (A) &:= \P(\omega \in \Omega: X(\omega) \in A)\\ 
	      &= \P( \{ (\omega_1, \omega_2) \in \Z \times \Z: \omega_1+ \omega_2 \in A \} \\
\end{align*}
Since $ A $ is a subset of $ \Z$ (which is countable) we have that $ A$ is also countable. This
means that the set $ \{ (\omega_1, \omega_2) \in \Z \times \Z: \omega_1+ \omega_2 \in A \} $ 
can be written as 
\[
    \{ (\omega_1, \omega_2) \in \Z \times \Z: \omega_1+ \omega_2 \in A \}  = \bigcup_{x_{i} \in A}
    \{ (\omega_1, \omega_2) \in \Z \times \Z: \omega_1+ \omega_2 = x_{i}  \} 
.\] 
Let's rewrite the set $ S = \{ (\omega_1, \omega_2) \in \Z \times \Z: \omega_1+ \omega_2 \in A \} $ one more time

\begin{align*}
   S &=  \bigcup_{x_{i}\in A}  \bigcup_{n \in \Z}  \{ (\omega_1, \omega_2) \in \Z \times \Z: \omega_1 = n, \, \omega_2 =
	 x_{i} -n \} \\
     &= \bigcup_{n \in \Z} \bigcup_{x_{i} \in A } (\{ (\omega_1, \omega_2) \in \Z \times \Z: \omega_1 = n, \, \omega_2 =
	 x_{i} -n \}) \\
\end{align*}
and the sets $ \bigcup_{x_{i} \in A} \{ (\omega_{1} , \omega_2): \omega_1= n, \omega_2 = x_{i} -n
\}$ are disjoint for different $ n$ since $ \omega_{1} = n $ is different
    for two sets with different $ n $. I am not sure how to prove that 
    we can change the order for two countable unions but in the finite case it is clear
    that it is possible to change the order of two finite unions.


We can apply the disjoint countable union property of probability distribution to get that 


\begin{align*}
    \P_{X}(A) &= \sum_{n \in \Z} \P( \bigcup_{x_{i} \in A} \{ (\omega_1, \omega_2): \omega_1 = n,
    \omega_2 = x_{i} -n \} )\\
\end{align*}
Similarly, the sets $ \{   (\omega_1, \omega_2): \omega_1 = n,
    \omega_2 = x_{i} -n\} $ are disjoint for different $ x_{i} $ since $ \omega_{2} = x_{i} -n $ is different
    for two sets with different $ x_{i} $s.

\begin{align*}
    \P_{X}(A) &= \sum_{n \in \Z} \sum_{x_{i}  \in A} \P( \{ (\omega_1, \omega_2): \omega_1 = n,
    \omega_2 = x_{i} -n \} )\\
	      &= \sum_{n \in \Z} \sum_{x_{i} \in A} p_{n}\cdot p_{x_{i} -n} \\
	      &= \sum_{x \in A} \sum_{n \in \Z} p_{n} p_{x-n} 
\end{align*}
where $ \P( \{ (\omega_1, \omega_2): \omega_1 = n,\omega_2 = x_{i} -n \} ) = p_{n} \cdot p_{x_{i}-n } $
by the definition of the probability distribution $ \P$.

We see that $ \P_{X} $ is a discrete distribution whose probability mass function is $ \sum_{n\in \Z} p_{n} p_{x-n}
$.
